\section{Proof of the main theorem}\label{sec:conclusion}

We first combine the three uniform bounds with the kernel energy estimate
for countable groups. We then pass through a finitely generated subgroup
to arbitrary discrete groups and include the classical averaging argument
for the reverse direction.

\begin{proof}[The countable nonamenable case]
Fix a countable discrete nonamenable group \(G\), and choose
\(S,d,\rho\) from \Cref{lem:contracting-average}.
For each \(\ell\ge1\), take the product list and assignment of
\Cref{prop:sparse-splitting}, with
\[
 n_\ell=d^\ell,\qquad
 r_\ell=\lceil n_\ell\rho^\ell\rceil,\qquad
 k_\ell=\lceil2r_\ell\log(2n_\ell)\rceil.
\]
In particular,
\begin{equation}\label{eq:dimension-ratio}
 \frac{k_\ell}{n_\ell}
 \le2\left(\rho^\ell+\frac1{n_\ell}\right)\log(2n_\ell)
       +\frac1{n_\ell}\longrightarrow0.
\end{equation}
Choose the vectors from \Cref{lem:frames} and let
\(P_i=v_iv_i^*\) at this stage. Use them to form the operators
\(A_\ell,T_\ell,U_\ell\) of \Cref{sec:weights} on
\(\mathcal H_\ell=\ell^2(G;\C^{k_\ell})\).
The row and column bounds of \Cref{cor:block-rows-columns} and the
commutator bound of \Cref{prop:uniform-differences} have the common
constant \(K=b_0^2=100\).

The invariant kernel of \(A_\ell\) is
\[
 W_\ell(s)=\sum_{i:s_i=s}P_i.
\]
It is finitely supported. Positivity of the weights gives
\begin{equation}\label{eq:kernel-energy}
 \sum_{s\in G}\norm{W_\ell(s)}_{\HS}^2
 =\sum_{i,j:s_i=s_j}\Tr(P_iP_j)
 =\sum_{i,j:s_i=s_j}|v_i^*v_j|^2
 \ge n_\ell.
\end{equation}
The diagonal terms equal one, and every other term is nonnegative;
thus collisions among group products only increase this lower bound.
By \Cref{eq:dimension-ratio}, the kernel energy divided by
\(k_\ell\) tends to infinity.

Fix \(\varepsilon>0\) and put \(t=\varepsilon/100\). Apply
\Cref{prop:direct-sum-obstruction} to the rescaled operators
\(tA_\ell,tT_\ell\). Their three uniform bounds have the common
constant \(tK=\varepsilon\). Their invariant kernels are
\(tW_\ell\), so their energies divided by \(k_\ell\) still tend to
infinity: they are \(t^2\) times the ratios just established.
The proposition therefore gives a representation on
\[
 \mathcal K=
 \left(\bigoplus_{\ell\ge1}\mathcal H_\ell\right)
 \oplus\left(\bigoplus_{\ell\ge1}\mathcal H_\ell\right)
\]
with uniform bound at most \(1+\varepsilon\) that is not similar to a
unitary representation. This Hilbert space is separable because \(G\)
is countable and each \(k_\ell\) is finite. This proves the nonamenable
direction and the quantitative assertion for countable groups.
\end{proof}

We next recall the induction argument for subgroup permanence of
unitarizability \cite[Proposition~0.5]{Pisier2004}, keeping track of its
norm and of the invariant inner product on the subgroup representation.

\begin{lemma}[Induction from a subgroup]\label{lem:subgroup-induction}
Let \(H\) be a subgroup of a discrete group \(G\), and let
\(\sigma:H\to\mathrm{GL}(\mathcal K)\) be a uniformly bounded
representation on a complex Hilbert space. There is a uniformly bounded
representation \(\Pi:G\to\mathrm{GL}(\mathcal L)\) on
\(\mathcal L=\ell^2(G/H;\mathcal K)\) such that
\[
 |\Pi|=|\sigma|.
\]
If \(\Pi\) is unitarizable, then \(\sigma\) is unitarizable.
\end{lemma}

\begin{proof}
Write \(X=G/H\) for the left cosets, and choose representatives
\(s:X\to G\) with \(s(H)=e\). For \(g\in G\) and \(x\in X\), put
\[
 c(g,x)=s(gx)^{-1}g s(x)\in H.
\]
The membership follows because \(s(gx)H=gs(x)H\). Direct multiplication
gives
\[
 c(gh,x)=c(g,hx)c(h,x)\qquad(g,h\in G,\ x\in X).
\]
On the dense subspace of finitely supported families in the Hilbert
direct sum \(\mathcal L=\bigoplus_{x\in X}\mathcal K\), define
\[
 (\Pi(g)\xi)_{gx}=\sigma(c(g,x))\xi_x\qquad(x\in X).
\]
If \(M=|\sigma|\), then reindexing by the permutation \(x\mapsto gx\)
gives
\[
 \norm{\Pi(g)\xi}^2
 =\sum_{x\in X}\norm{\sigma(c(g,x))\xi_x}^2
 \le M^2\sum_{x\in X}\norm{\xi_x}^2.
\]
Thus \(\Pi(g)\) extends to a bounded operator of norm at most \(M\) on
the whole direct sum, without a countability assumption on \(X\).
The cocycle identity gives \(\Pi(g)\Pi(h)=\Pi(gh)\) first on finitely
supported families and hence everywhere. Also \(\Pi(e)=I\), so
\(\Pi(g)^{-1}=\Pi(g^{-1})\). This proves \(|\Pi|\le M\).

Let \(j:\mathcal K\to\mathcal L\) be the isometric inclusion at the
identity coset \(H\). For \(h\in H\), one has \(hH=H\) and
\(c(h,H)=h\), whence
\[
 \Pi(h)j=j\sigma(h).
\]
In particular, \(\norm{\Pi(h)}\ge\norm{\sigma(h)}\) for every
\(h\in H\). Taking suprema gives the reverse inequality
\(|\Pi|\ge|\sigma|\), proving equality of the uniform bounds.

Suppose now that \([\cdot,\cdot]\) is an invariant equivalent inner
product for \(\Pi\). Let \(Q\) be the bounded positive invertible operator
representing this form. It satisfies
\[
 aI\le Q\le bI,\qquad \Pi(g)^*Q\Pi(g)=Q\quad(g\in G)
\]
for some \(0<a\le b<\infty\). The compression \(Q_H=j^*Qj\) satisfies
\[
 aI\le Q_H\le bI,\qquad
 \sigma(h)^*Q_H\sigma(h)
 =j^*\Pi(h)^*Q\Pi(h)j=Q_H\quad(h\in H).
\]
Consequently \([v,w]_H=[jv,jw]\) is an equivalent
\(\sigma(H)\)-invariant inner product on \(\mathcal K\). The compression
does not require \(Q\) to preserve the identity-coset summand. Hence
\(\sigma\) is unitarizable.
\end{proof}

\begin{proof}[Arbitrary discrete nonamenable groups]
Let \(G\) be a discrete nonamenable group and fix \(\varepsilon>0\).
The F\o lner criterion for an arbitrary discrete group says that
amenability is equivalent to the following condition: for every finite
\(S\subseteq G\) and every \(\delta>0\), there is a finite nonempty
\(F\subseteq G\) such that
\begin{equation}\label{eq:folner-finite-test}
 |sF\mathbin\triangle F|<\delta|F|\qquad(s\in S).
\end{equation}
This is the symmetric-difference form of F\o lner's Main Theorem
\cite[p.~245]{Folner1955}: inversion exchanges right and left
translations, and \(|sF\mathbin\triangle F|=2(|F|-|sF\cap F|)\).
No F\o lner sequence is required.

Since \(G\) is nonamenable, some finite \(S\) and positive \(\delta\)
fail this condition. The subgroup \(H=\langle S\rangle\) is finitely
generated and therefore countable. It is nonamenable: otherwise the
criterion applied inside \(H\) would give a finite nonempty
\(F\subseteq H\) satisfying \Cref{eq:folner-finite-test}, and the same
inequalities would hold in \(G\), a contradiction.

The countable case proved above gives a nonunitarizable
representation \(\sigma\) of \(H\) on a separable Hilbert space with
\(|\sigma|\le1+\varepsilon\). By Lemma~\ref{lem:subgroup-induction}, its
induced representation \(\Pi\) of \(G\) has
\(|\Pi|=|\sigma|\le1+\varepsilon\) and is nonunitarizable. This proves
the nonamenable direction for every discrete group.
\end{proof}

\begin{proof}[The amenable direction]
Let \(G\) be amenable and let \(\pi:G\to\mathrm{GL}(H)\) be uniformly
bounded, with \(M=|\pi|\). The Day--Dixmier argument
\cite{Day1950,Dixmier1950} works on an arbitrary complex Hilbert space.
If \(m\) is a left invariant mean, then
\(m_{\mathrm r}(f)=m(x\mapsto f(x^{-1}))\) is a right invariant mean.
Define a sesquilinear form by
\[
 [\xi,\eta]=m_{\mathrm r}\bigl(g\mapsto
                 \langle\pi(g)\xi,\pi(g)\eta\rangle\bigr).
\]
Since both \(\norm{\pi(g)}\) and \(\norm{\pi(g)^{-1}}\) are at most
\(M\), positivity of the mean gives
\[
 M^{-2}\norm{\xi}^2\le[\xi,\xi]\le M^2\norm{\xi}^2.
\]
Thus this is an equivalent Hilbert inner product. Right invariance
gives \([\pi(h)\xi,\pi(h)\eta]=[\xi,\eta]\) for every \(h\in G\).
Representing the form by a bounded positive invertible operator \(Q\),
we have \(\pi(h)^*Q\pi(h)=Q\). Hence \(S=Q^{1/2}\) is bounded and
invertible, and \(S\pi(h)S^{-1}\) is unitary for all \(h\). This proves
unitarizability and completes the equivalence.
\end{proof}

For uncountable nonamenable \(G\), the finitely generated subgroup \(H\) above has
uncountable index, and the induced witness is nonseparable. The conclusion
therefore uses the usual definition quantifying over all Hilbert spaces.
The argument concerns discrete groups and does not extend the theorem to
nondiscrete locally compact groups, where representations are normally
required to be strongly continuous.

\section{A consequence for group \texorpdfstring{$C^*$}{C*}-algebras}
\label{sec:cstar-consequence}

The following consequence uses the companion similarity theorem in
addition to \Cref{thm:main}.

Let \(G\) be any discrete group and let
\(\pi:G\to\mathrm{GL}(H)\) be a uniformly bounded nonunitarizable
representation on a complex Hilbert space. The operator-norm convergent sum
\[
 T_\pi(f)=\sum_{g\in G}f(g)\pi(g),\qquad f\in\ell^1(G),
\]
defines a bounded complex-linear unital homomorphism for convolution, with
\(\norm{T_\pi(f)}\le|\pi|\norm{f}_1\); absolute summability justifies
multiplication of the sums. Nevertheless, its restriction to \(\C[G]\)
has no bounded complex-linear unital algebra-homomorphism extension to
either the full group algebra \(C^*(G)\) or the reduced group algebra
\(C_r^*(G)\). Indeed, a full extension \(\Phi\) would be similar to a
\(*\)-homomorphism by the companion similarity theorem
\cite[Theorem~1.1]{OpenAIKadison2026}: some bounded invertible \(S\)
would make \(S\Phi(\cdot)S^{-1}\) a unital \(*\)-homomorphism.
For each canonical group unitary \(u_g\), one has \(\Phi(u_g)=\pi(g)\),
so every \(S\pi(g)S^{-1}\) would be unitary, a contradiction.
A reduced extension would compose with the canonical unital
\(*\)-quotient \(C^*(G)\to C_r^*(G)\) to give a full extension.
Thus the automatic \(\ell^1\)-norm bound does not imply
\(C^*\)-norm continuity. Applying this observation to the witnesses in
\Cref{thm:main} gives such a failure for every nonamenable discrete
group, even with uniform representation bounds arbitrarily close to one.

